%! Author = ewan
%! Date = 11/25/25

% Preamble
\documentclass[11pt]{article}

% Packages
\usepackage{amsmath}
\usepackage{xcolor}
\usepackage{hyperref}

% Document
% Document
\begin{document}

    \title{
            { \textbf{Security Protocols and Verification}} \\[1ex]
        {\small Attack of Cryptographic Protocols}
    }


    \author{
        Garance Frolla \\
        Ely Marthouret \\
        Ewan Decima\\ \\
        Team: \textbf{ASKO OM8464A2}
    }

    \date{September / November 2025}


    \maketitle
    \tableofcontents
    \newpage



    \section{Attack on Chronos V7}
    \subsection{Assumptions}
    We assume that the intruder $I$ initial knowledge $\mathcal{K}_I$ is
    \begin{center}
        $\mathcal{K}_I = \Bigl\{ K(A), K(B), K\bigl(\{|A|\}_{K_{BS}}\bigr), K(F) \Bigr\}$
    \end{center}

    Indeed, if $I$ plays the man-in-the-middle role and read all the messages in a legitimate conversation, $I$ is able to
    extract those information.

It is assumed that $I$ knows the function $F$, because if $F$ were kept secret between $B$ and $S$, the way it is shared would become questionable.
    In other words, if $F$ needs to remain secret between $B$ and $S$, it would have to be exchanged in the same way as a session key.





    \subsection{Scenario}
    
    \begin{enumerate}
        \item $A \rightarrow S : A, B, \{| K, N_K, A, B |\}$
        \item $S \rightarrow \textcolor{red}{I(B)} : S, \{|N_K|\}_{K_{BS}}, \{|K|\}_{N_K}, \{|A|\}_{K_{BS}}, \{|N_S|\}_{K_{BS}}$
        \item $I \rightarrow S : I, B, \{| \Sigma,  {N_I}_{\Sigma}, I, B |\}$
        \item $S \rightarrow \textcolor{red}{I(B)} : S, \{|{N_I}_{\Sigma}|\}_{K_{BS}}, \{|\Sigma|\}_{{N_I}_{\Sigma}}, \{|I|\}_{K_{BS}}, \{|N_S'\}_{K_{BS}}$
        \item $\textcolor{red}{I(S)} \rightarrow B : S, \{|{N_I}_{\Sigma}|\}_{K_{BS}}, \{|\Sigma|\}_{N_K}, \{|A|\}_{K_{BS}}, \{|N_S'|\}_{K_{BS}}$
        \item $B \rightarrow S : B, \{|N_S'|\}_{{N_I}_{\Sigma}}, \{|N_B|\}_{K_{BS}}$
        \item $S \rightarrow B : B, \{|F(N_S')|\}_{{N_I}_{\Sigma}}, \{|N_B|\}_{K_{BS}}$
        \item $B \rightarrow \textcolor{red}{I(A)} : B, \{|N_B|\}_{{N_I}_{\Sigma}}, \{|A|\}_{\Sigma}$
        \item $\textcolor{red}{I(A)} \rightarrow B  : A, \{|N_B|\}_{{N_I}_{\Sigma}}, \{|A|\}_{\Sigma}, \{|B|\}_{\Sigma}$

    \end{enumerate}

    \section{Attack Description}






    \subsection{Flow}

    \begin{itemize}
        \item \textbf{Message 1:} $A$ initiates the protocol, wanting to establish a secure session with $B$ through
        the trusted server $S$.

        \item \textbf{Message 2:} $S$ sends a legitimate message to $B$, but $I$ intercepts and blocks this communication.


        \item \textbf{Message 3:} $I$ starts a legitimate communication with $B$ involving the trusted server $S$ in
         order to get $\{{N_I}_{\Sigma}\}_{K_{BS}}$.




        \item \textbf{Message 4:} $S$ sends the legitimate message to $B$ (of $I$ - $B$ communication), but $I$ intercepts and blocks this communication.

        \item \textbf{Message 5:} $I$, pretending to be the $S$, sends $B$ the continuation of the communication involving $A$.
        thus choosing the key $\Sigma$, the associated nonce $\{|{N_I}_{\Sigma}\}_{K_{BS}}$, and replacing its encrypted identity with that of A

        \item \textbf{Message 6 and 7:} $B$ continues the run normally without the legitimate data


        \item \textbf{Message 8:} $I$ intercepts and blocks the message from $B$ intended for Alice. At this step, $I$ knows this
        \begin{center}
            $\mathcal{K}_I = \Bigl\{ K(A), K(B), K\bigl(\{|A|\}_{K_{BS}}\bigr), K(F),  K\bigl(\{|N_B|\}_{{N_I}_{\Sigma}}\bigr) \Bigr\}$
        \end{center}

        \item \textbf{Message 9:} $I$, pretending to be $A$, sends the final response to $B$ using the previous step.



    \end{itemize}

    \subsection{Result}

    In the end, $B$ believes he is communicating with $A$, but he is actually communicating with $I$ with a key chosen by $I$ : authentication fails.




\end{document}